% Section 13 — Conclusion (v4)

\section{Conclusion}
\label{sec:conclusion}

We argued that LLM agents do not perceive their own time. The argument has a theoretical layer and an empirical layer, and they agree.

The theoretical layer rests on two theorems. CIT (Theorem~\ref{thm:cit}) says no token-only loss can install chronoception. The Reverse-Scaling Theorem (Theorem~\ref{thm:reverse-scaling}) says that one of the most popular scaling strategies of the last two years --- reasoning-token expansion --- structurally degrades chronoception on the narrative axis where capability scaling otherwise helps. The two theorems together say: scale will not close the gap, and more reasoning will widen it.

The empirical layer measures ten agent configurations from three providers on roughly four thousand trajectories. Every panel agent honours less than five percent of the wall-clock budget given to it. Median self-narration error closes ninety-four percent along the non-reasoning capability axis but $\CAR$ does not change. The Reverse-Scaling Theorem is confirmed intra-model and cross-model with sign-flipping. Nominally ninety-percent confidence intervals on self-duration achieve zero to fifty percent coverage; the catastrophe is the calibration sub-problem's projection of CIT. No panel agent crosses the Augustine threshold $\ccestar = 0.20$, and the action axis alone exceeds it on every agent --- narrative training cannot close the gap. The three competing closed labs have, independently, installed wall-clock injection at the consumer chat tier and not at the developer-tool tier: converging engineering choices across competitors that carry the evidential weight of a natural experiment.

The framework's bridge from measurement to deployment is the Agentic Timeline Hypothesis. Chronoception is the binding constraint on autonomous-agent deployment horizons; the bound is bound by neither intelligence nor reasoning compute. Pre-registered P12 is qualitatively supported by 14{,}709 public METR HCAST runs --- reasoning models decay twenty percent faster than non-reasoning models at matched short-horizon success. The framework generalises along its natural other axis to the Cartographic Problem (Theorem~\ref{thm:sit}), the Agentic Frontier $T_{\max} \cdot S_{\max} \leq C/\ccestST$, and a planned third paper on the long-horizon agent benchmarks where spatial cognition is load-bearing.

We release \textsc{ChronoBench}, the Injection Atlas, the roughly four thousand trajectories, a Docker reproducibility package, a locked OSF pre-registration, and a \texttt{pip}-installable evaluation package. The Augustine threshold is uncrossable through narrative training; reasoning-token expansion makes it strictly harder; the spatial axis predicts the same structure. Closing it is the constructive program of \textsc{ChronoStack} (Paper 2) and \textsc{ChronoStack}$^{+}$ (Paper 3).

The systems currently called LLM agents do not perceive their own time. Until they do, they are tools we deploy for variable durations, not agents that inhabit time. We have given the gap a name, an ontology, two theorems, a benchmark, an atlas, a deployment bound, and an empirical confirmation across ten agents and four thousand trajectories. The next paper says how to close it.
